\scanpage{163}
\margin{162}
\begin{center}
INSTITUT DES HAUTES ETUDES SCIENTIFIQUES

Séminaire sur la topologie étale 1963-1964

Exposé VII

\uline{Topologies sur la catégories des préschémas.- Caractérisation}

\uline{``axiomatique'' de la topologie étale.}

par

JAN-ERIK ROOS
\end{center}

Nous allons maintenant introduire et étudier certaines topologies utiles sur $(\mathrm{Sch})=$ la catégorie des préschémas dans un univers $\underline{U}$ fixé dès maintenant (*). Nous utiliserons librement la terminologie et les résultats des exposés précédents. Souvent nous désignerons de la m\inserted{ê}me lettre un préschéma et son espace topologique sous-jacent. La m\inserted{ê}me remarque vaut pour un site et sa catégorie sous-jacente, une catégorie et l'ensemble de ses objets etc.

§ 1. \uline{La topologie de Zariski sur} $(\mathrm{Sch})$. \uline{Relations avec la géométri\inserted{e} algébrique ``classique''.}

Soit $X\in\underline{U}$ un préschéma. L'espace topologique sous-jacent à $X$ définit un site $\mathrm{Ouv}(X)\inserted{\in\underline{U}}$ (le site ``de Zariski'' de $X$ ). C'est la cohomologie de ce site (pour certains faisceaux et certains $X$) qui a été étudiée en détail dans la géométrie al\inserted{g}ébrique ``classique''. A vrai dire ce n'est pas exactement ce site qu'on a considéré au début, mais seulement le site analogue défini sur le sous-espace $X_o$ de $X$ formé des points \uline{fermés} de $X$. Précisons ceci: On dit qu'un préschéma $X$ est un \uline{préschéma de Jacobson} si le foncteur $\mathrm{Ouv}(X)\xrightarrow{\cap}\mathrm{Ouv}(X_o)$ défini par $U\rightsquigarrow U\cap X_o$ induit une équivalence des catégories des faisceaux d'ensembles (ou une équivalence des catégories des faisceaux abéliens, c'est la m\inserted{ê}me chose $\ldots$ ) sur ces deux sites.

\margin{\figure{\begin{tikzpicture}\draw (0,0)--(5.5,0);\end{tikzpicture}} (*) Nous dirons qu'un préschéma $(X,O_X)$ appartient à l'univers $\underline{U}$ si $X\in\underline{U}$ et si de plus $O_X$ est un faisceau de $\underline{U}$-ensembles.}
